\documentclass[12pt,ifthen]{article}
\usepackage{url}
\usepackage{comment}

\newif{\ifshowsoln}
% \showsolntrue

\newcommand{\und}{\_\_\_\_\_\_\_\_\_}
\newcommand{\Z}{\mathbb{Z}}
\usepackage{amsmath}
\usepackage{amssymb}  % for \nmid

\begin{document}

\centerline{\bf Gurkeerat Singh}
\vspace{0.3cm}
\centerline{\bf HW 00 CMSC 250}
\centerline{\bf Morally Due Feb 3, 2026,  9:30AM} 
\centerline{\bf Dead Cat Due Feb 5, 2026,  9:30AM} 

\bigskip

\centerline{\bf This is the LAST HW where I will}
\centerline{\bf Put the Dead Cat Reminder.}
\centerline{\bf In the future ``Morally Due'' will imply Dead Cat.}


\bigskip

\noindent
{\bf Course Website}

\url{https://www.cs.umd.edu/~gasarch/COURSES/250/S26/index.html}

\begin{enumerate}
\item
(24 points- 3 points each) READ the Syllabus before doing this problem. 
You may also need to look at the slides from Jan 27. 
\begin{enumerate}
\item 
When is the Midterm? {\bf The midterm is on March 31, 2026.}
\item
What is {\it the dead-cat policy}? {\bf While the deadline for homework is at 9:30 A.M. on Tuesdays, everyone gets an extension to Thursday morning in case a student is truly unable to complete it in time due to some other priority.}
\item
Why is is called {\it the dead-cat policy}? {\bf This extension isn't to be treated like the real due-date. Rather, it's reserved for sudden and unexpected situations that take precedence over completing homework, like the death of a cat.}
\item
If I write that a HW is {\it Morally Due Feb 25, 2026} what does that mean? \textbf{\textit{Morally Due Feb 25, 2026} means that the due date is to be treated like Feb 25, but the Dead-Cat Policy will allow for submissions up to 9:30 A.M on Feb 27, 2026.}
\item
When are Leo's Office Hours? {\bf Leo's Office Hours are from 2-4 P.M. on Mondays at IRB1266, or by appointments on Zoom.}
\item
When are Bill's Office Hours? {\bf Bill's Office Hours are from 11 A.M-12 P.M, and from 2-3 P.M in IRB2242, or also by appointments over Zoom.}
\item
Name all of the ways you can communicate with Leo and Bill. {\bf You can communicate with Leo and Bill during lecture, discussion, office hours, email, Piazza, or via a Zoom Appointment.}
\item
Describe the In-Class assignments. How many can you miss without
penalty. {\bf The In-Class assignments will make the students get into groups to work on a problem, and they will be at random. There will be around 10 during the semester and each one will be worth 1\% of your final grade. You can miss up to 3 of these without penalty given that students sometimes may not be in class.}
\end{enumerate}


\newpage


\item
(36 points-12 each)
Learn Python 3 and write a simple program in it before doing this problem. 
\begin{enumerate}
\item
Name a PRO of Python. {\bf A pro of Python is that it's very similarly to pseudocode and almost English, which makes it very beginner-friendly and easier to understand.}
\item
Name a CON of Python. {\bf Since Python is an interpreted langauge, one of its biggest cons is that it tends to be slower than other languages like C at runtime.}
\item
What did your program do? {\bf I created a Blackjack Simulator which gave the dealer one card face up and gave me 2 at random. Based on basic Blackjack strategy, the program uses the values of the cards to decide whether to hit, or stand. After the player is done hitting, the dealer gives cards to themselves after they exceed 16. Finally, the sums of the cards are compared to determine the winner of the game. I enjoyed tweaking the algorithm of the program and running it a large amount of times to optimize our program's win percentage.}
\end{enumerate}

\newpage 

\item
(0 points but do it anyway. This is the most important problem on this HW!)
\begin{enumerate}
\item 
Learn LaTeX and write some simple documents in it. 
\item 
Write a LaTeX document that summarizes the lecture on Propositional Logic. 
The document should include a truth table (so you get practice doing
LaTeX tables) and some formulas on the variables $x,y,z$
(so you get practice typesetting mathematics). 

\begin{center}
    \begin{tabular}{|c|c|c|}
        \hline
        p & q & $p \wedge (\sim q)$ \\
        \hline
        {\bf F} & {\bf F} & {\bf F} \\
        \hline
        {\bf F} & {\bf T} & {\bf F} \\
        \hline
        {\bf T} & {\bf F} & {\bf T} \\
        \hline
        {\bf T} & {\bf T} & {\bf F} \\
        \hline
    \end{tabular}
\end{center}

\bf A few laws using the variables \it x, y, and z :

\bf Associative Law:

$ (x \vee y) \vee z \equiv x \vee (y \vee z) $

$ (x \wedge y) \wedge z \wedge x \wedge (y \wedge z) $

\bf Law of Absorption: 

$ x \vee (y \wedge (x \vee z)) \equiv x \vee (y \wedge z) $

\bf Distributive Law:

$ x \wedge (y \vee z) \equiv (x \wedge y) \vee (x \wedge z) $
\end{enumerate}


\newpage

\item
(40 points) 
Read Vannevar Bush's paper from July 1945: 

\url{http://web.mit.edu/STS.035/www/PDFs/think.pdf}

Write down two predictions he made that came true. 

(Note- this is not a paper in Discrete Math but it is such a good paper 
that every undergraduate should read it!)

\bf Two predictions Vannevar Bush made in his 1945 paper that came true were the ability to store millions of volumes digitally, or a digital library. This preceded today's technologies like Google Books and cloud storage systems that allow instant access to a large amount of data. Bush also predicted the future of speech-to-text technology, where machines could transform spoken words into written text. We see this concept in common everyday applications like Siri, Google Assistant, and Alexa.  

\end{enumerate}
\end{document}

